% ============================================================
% common/questions_env.tex —— 题目 / 解析的语义包装环境
%
% questions.tex 中每道题包在 \begin{question}...\end{question} 里，
% answers.tex 中每题解析包在 \begin{answer}...\end{answer} 里；
% 题头排版仍由内容行中的 \gkti / \gksolhead 完成（见 macros.tex）。
%
% 环境本身不产生任何盒子或垂直间距（只引入一个 TeX 分组），
% 排版结果与不包裹时逐点一致。它的价值在于：
%   1. 给每道题一个机器可数的结构边界（校验脚本数环境数 = 题量）；
%   2. 为将来的条件编译留钩子——某版本可重定义环境实现按题取舍。
% ============================================================

% 环境不把整题装盒：题目仍可自然跨页，只在题头处要求保留极少空间，
% 仅避免题号/解析标题恰好成为页尾孤行（第一行与后续内容完全分离）。
% 保留量必须小：\Needspace{8\baselineskip} 会把最后一题整题推到新页，
% 在上一页底留下大面积空白，反而制造孤页（见历史卷 review 的
% pagination_whitespace fail）。2 个 baseline 足以让题头 + 首行同行，
% 又不会在「只剩几行」时强制换页。
\newenvironment{question}{\Needspace{2\baselineskip}}{}
\newenvironment{answer}{\Needspace{2\baselineskip}}{}
